\section{Plan Mode：通过协作生成规格}

Codex 的第一种规划方式是 Plan Mode。用户可以通过 \texttt{/plan} 进入这个模式。进入之后，系统会在 developer message 中植入一组完整的行为约束，让模型先探索、再澄清、再形成一份可执行的 Markdown 计划。

\begin{figure}[H]
\centering
\includegraphics[width=\textwidth]{frames/frame-00-03-35.jpg}
\caption{Plan Mode 的核心不是立刻执行，而是通过对话形成可实施规格。\protect\footnotemark}
\end{figure}
\footnotetext{视频画面时间区间：00:03:25--00:03:45。}

Plan Mode 与一个特殊工具配合：\texttt{request\_user\_input}。这个工具只在 Plan Mode 中可用，用来向用户提出一到三个短问题，等待用户回答，从而减少需求不清导致的错误实现。

\begin{knowledgebox}{Plan Mode 的产物在哪里？}
讲者特别强调：Plan Mode 结束后形成的是一段 Markdown 文本，但它并不会自动落盘成本地文件。它存在于 context window 中，作为后续执行阶段的上下文。换句话说，这是一份上下文内的规格，而不是仓库里的持久化文档。
\end{knowledgebox}

Plan Mode 的约束也很强：在这个阶段应该做非变更式探索，不直接修改代码。它更像是一次需求澄清和规格设计会话，最后再询问用户是否执行这份计划。用户确认后，系统回到默认执行模式。

\subsection{Plan Mode 的适用场景}

Plan Mode 适合需求边界还不清楚、实现路径有多个分叉、或者需要用户选择取舍的任务。例如一个大型重构可能涉及 API 兼容性、数据迁移和测试策略，此时先让模型用问题收敛规格，比直接动手更稳。

\begin{warningbox}{Plan Mode 不是持久化项目管理工具}
Plan Mode 的计划只是上下文对象。它可以指导当前会话，但如果希望跨会话保存、恢复或审计，仍然应该像本仓库一样使用文件化计划，例如 \texttt{task\_plan.md}、\texttt{NOTE\_GENERATION\_TODO.md} 等。
\end{warningbox}

\subsection{本章小结}

Plan Mode 是一种协作模式。它通过 developer instruction 和 \texttt{request\_user\_input} 约束模型行为，把``先问清楚再写代码''产品化成一个明确流程。

